\section {13 rocks (Number theory)}

You are given 13 rocks. If you remove any one of them, the remaining 12 can be split
into two groups of 6 rocks each that balance each other on a scale.
Prove that all 13 rocks have the same weight.

\subsection{Solution}

First assume that all weights $w_1, \ldots, w_{13}$ are whole numbers and let $S$ be their sum.
If we remove rock $i$, the remaining 12 rocks are split into two groups of equal weight, so $S-w_i$ is even.
This holds for every $i$, so all the $w_i$ have the same parity as $S$: they are either all even or all odd.

Notice that two operations keep the property of the rocks (any 12 balance as 6 against 6):
subtracting the same number from all weights (each side of the scale holds 6 rocks, so both sides lose the same amount)
and dividing all weights by the same number.

Now subtract the smallest weight from all weights. The property still holds and one of the weights is $0$.
Since all weights have the same parity as that $0$, they are all even.
Divide all of them by $2$. They are still whole numbers, the property still holds, so they are again all even, and we can divide by $2$ again.
We can repeat this forever, which is only possible if all the weights are $0$.
That means that before subtracting, all the weights were equal to the smallest weight.

If the weights are rational, multiply all of them by a common denominator. That gives whole numbers with the same property, which are therefore all equal.

If the weights are any real numbers the argument needs one more step.
Each balancing is a linear equation in $w_1, \ldots, w_{13}$ whose coefficients are $0$, $1$ and $-1$.
The argument above shows that the only rational solutions of this system of equations are of the form $(c, c, \ldots, c)$,
so the solution space over the rationals has dimension $1$.
The dimension of the solution space of a system of linear equations is $13$ minus the rank of the matrix, and the rank
is computed by Gaussian elimination, which never leaves the rationals when the coefficients are rational.
So the rank is the same over the reals and the solution space over the reals also has dimension $1$: all the weights are equal.
